\documentclass[12pt]{article}

\usepackage[portuguese]{babel}
\usepackage[utf8]{inputenc}
\usepackage{graphicx}
\usepackage{geometry}
\usepackage{hyperref}
\usepackage{fancyhdr}
\usepackage{setspace}
\geometry{a4paper, margin=2.5cm}
\setstretch{1.3}

\pagestyle{fancy}
\fancyhf{}
\rhead{Laboratórios de Informática}
\lhead{Dashboard PS2}
\rfoot{\thepage}

\title{
\Huge \textbf{Trabalho Prático}\\[10pt]
\Large Laboratórios de Informática\\[10pt]
\Large Dashboard PS2
}
\author{
\textbf{Autores:}\\
Xico Dias\\
(Nome do colega)
}
\date{2026}

\begin{document}

\maketitle
\newpage

\tableofcontents
\newpage

\section{Introdução}

O presente trabalho foi realizado no âmbito da disciplina de Laboratórios de Informática e teve como objetivo o desenvolvimento de uma aplicação em Python capaz de ler, validar, processar e analisar ficheiros PS2 (relacionados com Débitos Diretos), apresentando os resultados através de um dashboard interativo desenvolvido com a biblioteca Shiny para Python.

Este projeto permitiu trabalhar competências fundamentais como:
\begin{itemize}
\item Programação em Python
\item Manipulação de ficheiros estruturados
\item Processamento e análise de dados
\item Organização modular de código
\item Utilização de Git e GitHub
\end{itemize}

\section{Objetivos do Trabalho}

Os principais objetivos definidos foram:

\begin{itemize}
\item Ler e interpretar ficheiros PS2
\item Validar a estrutura e organização dos dados
\item Extrair informação relevante
\item Calcular estatísticas sobre os valores
\item Apresentar os resultados de forma clara e intuitiva
\item Criar um dashboard funcional e interativo
\item Utilizar ferramentas de desenvolvimento colaborativo
\end{itemize}

\section{Arquitetura do Sistema}

O projeto foi organizado de forma modular na pasta \texttt{app}, com três ficheiros principais:

\subsection{parser.py}

Responsável por:
\begin{itemize}
\item Leitura dos ficheiros PS2
\item Validação do formato
\item Extração de campos relevantes
\item Transformação para estrutura utilizável
\end{itemize}

\subsection{analyzer.py}

Responsável por:
\begin{itemize}
\item Construção da tabela de dados
\item Cálculo de estatísticas:
\begin{itemize}
\item Total recebido
\item Número de operações
\item Valor máximo
\item Valor mínimo
\item Média
\end{itemize}
\end{itemize}

\subsection{dashboard.py}

Responsável por:
\begin{itemize}
\item Interface gráfica web
\item Apresentação das estatísticas
\item Exibição da tabela
\item Geração do gráfico
\end{itemize}

\section{Dashboard Desenvolvido}

O dashboard apresenta:
\begin{itemize}
\item Cartões informativos com estatísticas principais
\item Tabela de operações processadas
\item Gráfico com análise visual dos valores
\end{itemize}

\section{Tecnologias Utilizadas}

\begin{itemize}
\item Python 3
\item Shiny para Python
\item Pandas
\item Matplotlib
\item Git e GitHub
\end{itemize}

\section{Resultados Obtidos}

O sistema final encontra-se:
\begin{itemize}
\item Funcional
\item Organizado
\item Modular
\item Testado com ficheiros PS2
\end{itemize}

O dashboard apresenta corretamente os dados, estatísticas e gráficos exigidos no enunciado.

\section{Conclusão}

Todos os objetivos propostos no trabalho foram cumpridos com sucesso.  
O projeto permitiu consolidar conhecimentos de programação, análise de dados e desenvolvimento de aplicações interativas, bem como a utilização de ferramentas profissionais de gestão de código.

\section{Repositório GitHub}

O projeto encontra-se disponível em:

\begin{center}
\textbf{\url{https://github.com/xicodias7-byte/dashboard-ps2}}
\end{center}

\end{document}
